
\begin{theorem}
    \textbf{Teorema del valor medio para integrales dobles.}   Sea $f:D\subseteq\Rn{2}\to\R$ continual.  Entonces existe $(x_0,y_0)$ en $D$ tal que
    \[
        \int_D f(x,y)\:dA=f(x_0,y_0)\text{A}(D).
    \]
\end{theorem}

\begin{definition}
    Sea $\mathbf{F}:A\subseteq\Rn{n}\to\Rn{n}$ un campo $\mathcal{C}^1(A)$ y sea $\boldsymbol{D}_\mathbf{F}$ (\autoref{def:mat_dif}) su matriz diferencial.  Se define el \emph{jacobiano}, o \emph{determinante jacobiano}, y se lo nota $J_\mathbf{F}$ a 
    \[
        J_\mathbf{F}=\text{det}(\boldsymbol{D}_\mathbf{F}).
    \]
\end{definition}


\begin{example}  Calcular $J_\mathbf{F}$ siendo  \( \mathbf{F}(r, \theta) = (r\cos\theta, r\sin\theta) \).

 Primero, calculamos su la matriz diferencial.
       \[
       \boldsymbol{D}_{\mathbf{F}}(r, \theta) 
       =
       \begin{bmatrix}
       \cos\theta & -r\sin\theta \\
       \sin\theta & r\cos\theta
       \end{bmatrix}
       \]
       Entonces, el jacobiano es
       \[
       J_{\mathbf{F}}(r, \theta) = \det(\boldsymbol{D}_{\mathbf{F}})(r, \theta) = 
      r(\cos^2\theta + \sin^2\theta) = r.
       \]
   \end{example}
   
   \begin{example}  Calcular $J_\mathbf{F}$   siendo   \( \mathbf{F}(r, \theta, z) = (r\cos\theta, r\sin\theta, z) \).
        
  Primero, calculamos su la matriz diferencial.
   \[
    \boldsymbol{D}_{\mathbf{F}} (r, \theta,  z)=
    \begin{bmatrix}
   \cos\theta & -r\sin\theta & 0 \\
   \sin\theta & r\cos\theta & 0 \\
   0 & 0 & 1
   \end{bmatrix}
   \]   
    Entonces, el jacobiano es 
   
   \[
   J_\mathbf{F} (r, \theta, z) = \det(D_F)(r, \theta, z)  = 
   \begin{vmatrix}
   \cos\theta & -r\sin\theta & 0 \\
   \sin\theta & r\cos\theta & 0 \\
   0 & 0 & 1
   \end{vmatrix}
   = r
   \]
   \end{example}
   
   \begin{example}  Calcular $J_\mathbf{F}$   siendo  \( \mathbf{F}(\rho, \theta, \phi)  = (\rho\sin\phi\cos\theta, \rho\sin\phi\sin\theta, \rho\cos\phi) \). 
   
   
    Primero, calculamos su la matriz diferencial.
   
    \[
    \boldsymbol{D}_{\mathbf{F}}(\rho, \theta, \phi) 
    =
    \begin{bmatrix}
    \sin\phi \cos\theta & -\rho \sin\phi \sin\theta & \rho \cos\phi \cos\theta \\
    \sin\phi \sin\theta & \rho \sin\phi \cos\theta & \rho \cos\phi \sin\theta \\
    \cos\phi & 0 & -\rho \sin\phi
    \end{bmatrix}
    \]
     Entonces, el jacobiano es 
    
    \[
    J_\mathbf{F}(\rho, \theta, \phi)  = \det(\boldsymbol{D}_{\mathbf{F}})(\rho, \theta, \phi)  = \begin{vmatrix}
        \sin\phi \cos\theta & -\rho \sin\phi \sin\theta & \rho \cos\phi \cos\theta \\
        \sin\phi \sin\theta & \rho \sin\phi \cos\theta & \rho \cos\phi \sin\theta \\
        \cos\phi & 0 & -\rho \sin\phi
        \end{vmatrix}=\rho^2 \sin\phi
    \]
    
    
    \end{example}
